\documentclass[11pt]{article}
\usepackage{amsmath,amssymb,amsthm,hyperref}
\newtheorem{theorem}{Theorem}
\newtheorem{lemma}[theorem]{Lemma}
\begin{document}
\title{An Equality Kernel in Paired Multilinear Coefficients}
\author{}
\date{November 24, 2026}
\maketitle

\paragraph{Status.}
This is a complete written proof of a coefficient programming mechanism,
followed by a counterexample to its proposed global assembly. None of
this note is Lean-checked. It proves no asymptotic bound for
\texttt{worstCaseHeadComplexity}. In particular, coefficients of the
multilinear interpolant are not the values of the rational score on the
Boolean cube. The companion note
\texttt{continuation-coefficient-bridge-audit.tex} gives an independent
audit of this distinction and a bounded incidence-span theorem.

\section{Two heads and an exact coefficient support restriction}

Partition a finite coordinate set into row pairs $(x_i,x'_i)$,
column pairs $(y_j,y'_j)$, and reserve coordinates $z_k$. Choose
positive weights $a_i,b_j,r_k$ and $\epsilon>0$, and write
\[
 D=\epsilon+\sum_i a_i(2-x_i-x'_i)
       +\sum_j b_j(2-y_j-y'_j)+\sum_k r_k(1-z_k),
 \quad A=\sum_i c_i(x_i-x'_i),\quad
 B=\sum_j d_j(y_j-y'_j).
\]
For every sufficiently small $t>0$, define
\[
 F_t=\frac1{2t}\left(\frac{A}{D-tB}-\frac{A}{D+tB}\right)
     =\frac{AB}{D^2-t^2B^2}.                         \tag{1}
\]
Every slope of both denominators is strictly negative for sufficiently
small $t$, and both denominators are strictly positive on the whole
cube. Indeed, their slopes tend to the strictly negative slopes of
$D$, while their values tend uniformly to values at least $\epsilon$.
Thus (1) consists of two original admissible affine-ratio heads.
On the finite cube it converges uniformly to $F_0=AB/D^2$.

Let $S$ be the full coordinate set. The unique multilinear interpolant
of $F_t$ has zero coefficient at $x^S$, at every monomial obtained by
deleting one coordinate, and at every monomial obtained by deleting two
coordinates except a row coordinate and a column coordinate. To see
this, interchange both coordinates in every row pair. This negates
$A$ and preserves $D,B$, hence negates $F_t$. Interchanging every
column pair similarly negates $F_t$. Every asserted zero monomial is
invariant under at least one of these two interchanges. Uniqueness of
the multilinear interpolant therefore forces its coefficient to be
its own negative.

\section{The positive Gram kernel}

Put
\[
 W(u)=\prod_i(1-e^{-a_i u})^2
             \prod_j(1-e^{-b_j u})^2
             \prod_k(1-e^{-r_k u}),\qquad
 q(v)=\frac{e^{-v}}{(1-e^{-v})^2},
\]
and define
\[
 K(a,b)=\int_0^\infty u e^{-\epsilon u}
                         W(u)q(au)q(bu)\,du.        \tag{2}
\]
The coefficient of $x^{S\setminus\{x_i,y_j\}}$ in $F_0$ is
\[
                         c_i d_j K(a_i,b_j).        \tag{3}
\]
Deleting a primed coordinate instead changes the sign for that pair.

For a direct verification, use $D^{-2}=\int_0^\infty u e^{-Du}du$.
An untouched pair of equal-weight coordinates contributes
$(1-e^{-au})^2$ to its full multilinear coefficient. After deleting
$x_i$, the contribution of $c_i(x_i-x'_i)$ and its pair is
$-c_i e^{-a_i u}$. Every other term of $A$ has zero full coefficient
on its own pair. Dividing the exceptional contribution by the
untouched pair factor gives $-c_iq(a_i u)$. The two minus signs from
the row and column deletions give (3). This computation also shows
that there is no alternating sign depending on $|S|$ in (2).

Equation (2) is a Gram kernel for the functions $q(au)$ in the
positive measure $u e^{-\epsilon u}W(u)du$. Its diagonal entries are
strictly positive and finite. Choosing
$c_i=K(a_i,a_i)^{-1/2}$ and $d_j=K(b_j,b_j)^{-1/2}$ normalizes (3)
to its correlation kernel.

\section{Geometric labels become orthogonal}

\begin{theorem}\label{thm:orthogonal}
Fix integers $L\ge1$ and $m_0,\ldots,m_L\ge2$. For $R>1$ put
$a_g=R^g$, and replace $W$ in (2) by
\[
                  W_R(u)=\prod_{\ell=0}^L
                                  (1-e^{-R^\ell u})^{m_\ell}.
\]
For every $g,b\in\{1,\ldots,L\}$,
\[
 \lim_{R\to\infty}
 \frac{K_R(R^g,R^b)}
 {\sqrt{K_R(R^g,R^g)K_R(R^b,R^b)}}
                       ={\bf1}_{g=b}.              \tag{4}
\]
The label $0$ is reserved and is not queried.
\end{theorem}

\begin{proof}
Define $M_g=\sum_{\ell<g}m_\ell$ and
\[
                  E_g=-2g+\sum_{\ell<g}m_\ell(\ell-g).
\]
Substitution $u=R^{-g}v$ gives
\[
 R^{-E_g}K_R(R^g,R^g)\longrightarrow
 C_g=\int_0^\infty
       v^{1+M_g}(1-e^{-v})^{m_g}q(v)^2\,dv>0.        \tag{5}
\]
For $g>b\ge1$, the same substitution gives
\[
 R^{-E_g-2(g-b)}K_R(R^g,R^b)\longrightarrow
 C_{gb}=\int_0^\infty
       v^{M_g-1}(1-e^{-v})^{m_g}q(v)\,dv>0.          \tag{6}
\]
Here $(1-e^{-\rho v})/\rho\to v$ and
$\rho^2q(\rho v)\to v^{-2}$ as $\rho\downarrow0$.
For domination, use $(1-e^{-\rho v})/\rho\le v$,
$q(v)\le v^{-2}$, and
$1-e^{-v}\le\min(1,v)$. At zero the resulting power in both
integrands is $v^{M_g+m_g-3}$, which is integrable because
$M_g\ge m_0\ge2$ and $m_g\ge2$. At infinity the factor $q(v)$
or $q(v)^2$ gives exponential decay. High-label factors are bounded
by one. These bounds justify (5) and (6) by dominated convergence.

It remains to show that the power in the normalized off-diagonal
entry is negative. For $g>b$ it is
\[
 T_{gb}=\frac{E_g-E_b}{2}+2(g-b).
\]
The recurrence
$E_{k+1}-E_k=-2-\sum_{\ell\le k}m_\ell$ gives
\[
 2T_{gb}=\sum_{k=b}^{g-1}
                  \left(2-\sum_{\ell\le k}m_\ell\right)<0,
\]
because $b\ge1$ and each sum is at least $m_0+m_1\ge4$.
Equations (5) and (6) prove (4). The diagonal case is identically one.
\end{proof}

The theorem permits arbitrary fixed multiplicities, so different
two-head modules need not use a common multiset of labels. First
choose $R$ to make all finitely many undesired correlations small;
then choose $t$ in (1) small enough to preserve all required coefficient
inequalities. The exact coefficient zeros established above survive
for every such $t$. This is a genuine equality mechanism in the
specified coefficient slice.

\section{Why this does not yet assemble a Boolean classifier}

Consider only one row pair and one column pair, both of weight $R$,
two reserve coordinates of weight one, and $\epsilon=1$. Take
$A=x-x'$ and $B=y-y'$. In $F_0=AB/D^2$, delete $x$ and $y$.
The coefficient including both reserves equals
\[
 K_4=\frac1{(2R+1)^2}-\frac2{(2R+2)^2}
                      +\frac1{(2R+3)^2}
                  \sim\frac3{8R^4}.
\]
The coefficient omitting one reserve, and the coefficient omitting
both, respectively equal
\[
 K_3=\frac1{(2R+2)^2}-\frac1{(2R+3)^2}
               \sim\frac1{4R^3},\qquad
 K_2=\frac1{(2R+3)^2}\sim\frac1{4R^2}.
\]
Thus $K_3/K_4\sim2R/3$ and $K_2/K_4\sim2R^2/3$.
Normalizing the intended coefficient to one makes lower coefficients
arbitrarily large. The mechanism does not have the small lower tails
needed for a descending-coefficient sign construction.

There is also no legitimate shortcut by evaluating the multilinear
interpolant at rescaled inputs. For example, the Boolean interpolant
of $1/(1+x)$ is $1-x/2$. At $x=3$ the interpolant is negative and
the rational expression is positive. Substitution into the rational
head does not perform substitution into its Boolean interpolant.

Finally, at actual mixed-corner vertices the simpler score
$AB/D^2$, with $A=\sum_i a_i(1-x_i)$ and
$B=\sum_j b_j(1-y_j)$, has normalized value
$4a_i b_j/(\epsilon+a_i+b_j)^2$. This is the already known forward
equality kernel in deletion coordinates. Adding a large positive
coefficient on coordinates outside a downset center produces mixed
denominator slopes. It does not supply the reverse-grid compiler
required by the gluing theorem.

The exact remaining issue is a whole-cube assembly preserving strict
output signs at the improved head count. Coefficient orthogonality,
bounded matrix decompositions, and small parameter $t$ do not solve
that issue. This route is closed unless a new assembly mechanism or
a genuinely new lower-tail invariant is supplied.
\end{document}
